\textbf{V -- ATTACCO E CONTRATTACCO}

La distinzione tra \emph{iniziativa} e \emph{attacco} negli scacchi è di
contenuto e di gradazione.

Uno dei due colori ha l'iniziativa quando è in grado di portare minacce
in posizioni prive di debolezze manifeste. Tali ripetute minacce servono
a indebolire la posizione avversaria a tal punto da permettere un
attacco vero e proprio; in tal senso l'iniziativa può essere vista come
preludio all'attacco.

L'attacco è il tentativo di sfruttare tali debolezze e mira a un
vantaggio decisivo.

In posizioni prive di debolezze le manovre dei giocatori tendono a
conquistare l'iniziativa e, con piccole e grandi minacce, a creare le
debolezze necessarie a passare alla fase di attacco vero e proprio.
Quando il centro delle manovre è lo stesso settore della scacchiera per
entrambi, allora uno solo dei due giocatori avrà l'iniziativa,
altrimenti un giocatore può avere l'iniziativa in un settore (per
esempio, l'ala di Re) mentre l'avversario l'avrà sull'altro lato della
scacchiera. In questi casi lo spostamento dei pezzi da un settore
all'altro può facilmente essere causa di indebolimenti e di scompensi;
resta il fatto che in mancanza di debolezze un attacco è destinato
all'insuccesso.

Questo è quasi sempre vero. Le eccezioni dipendono da particolari
strutture di pedoni oppure dall'ubicazione dei Re.

\bookdiagramplaceholder{assets/diagrams/image74.png}{74}

La posizione del diagramma, in mancanza di una rimozione del §e5,
suggerisce un'iniziativa bianca sul lato di Re, cui farà da contrapposto
quella nera ad ovest. Il Bianco opererà la spinta f2-f4-f5 mentre il
Nero quella c7-c5.

In caso di arrocchi ad est, e se il centro rimane chiuso, il Bianco
potrà anche permettersi di spingere i pedoni \emph{g} e \emph{h}
imbastendo un attacco in piena regola.

\bookdiagramplaceholder{assets/diagrams/image75.png}{75}

In quest'altro schema, invece, le parti sono invertite. Il Nero ha
l'iniziativa ad est e il Bianco sull'ala di Donna. Il Nero effettuerà la
spinta f7-f5, in caso di cambio può in seguito contrapporre
all'iniziativa Bianca ad ovest (b2-b4 e c4-c5) l'apertura del centro con
e5-e4, altrimenti, a centro chiuso, può giocare f5-f4 e attaccare sul
lato di Re, eventualmente anche avvalendosi delle spinte dei pedoni
\emph{g} e h.

La partita che segue, vinta dal Campione del Mondo Garry Kasparov con il
consueto slancio ricco di complicazioni, illustra bene le potenzialità
illustrate dalla struttura diagrammata.

Piket--Kasparov, Tilburg 1989 \textbf{(Difesa Est Indiana)}

\textbf{1.d4 ¤f6 2.¤f3 g6 3.c4 ¥g7 4.¤c3 0-0 5.e4 d6} nella scelta delle
difese a 1.d4 Kasparov ripercorre le vie di Fischer. Entrambi
prediligono l'Est Indiana e la Gruenfeld.

\textbf{6.¥e2 e5 7.0-0 ¤c6 8.d5 ¤e7 9.¤e1 ¤d7 10.¥e3 f5 11.f3 f4 12.¥f2
g5} come abbiamo detto, in presenza di un centro bloccato ci si può
permettere di spingere i pedoni davanti all'arrocco.

\textbf{13.b4 ¤f6 14.c5 ¤g6 15.¤xd6 ¤xd6 16.¦c1 ¦f7 17.a4 ¥f8 18.a5 ¥d7
19.¤b5 g4!} gli attacchi sulle ali sono in pieno svolgimento. Inferiore
al tratta del testo era 19... ¥xb5 20.¥xb5 g4 21.fxg4 ¤xe4 22.¥d3.

\textbf{20.¤c7?!}

\bookdiagramplaceholder{assets/diagrams/image76.png}{76}

Una mossa di dubbio valore. Il Bianco avrebbe fatto bene a giocare
20.fxg4! ¤xe4 21.¤c7 ¥a4!? \emph{(21... ¦c8 22.¤e6 £f6 23.¦xc8 ¥xc8
24.¥c4)}.

Anche la presa 20.¤xa7 sarebbe stata insufficiente: 20... g3 21.¥b6 £e7
\emph{(21... gxh2+? 22.¢xh2! £e7 23.¦h1! ch5 24.¢g1 ¤g3 25.¦h2)} 22.¤b5
\emph{(22.h3? ¥xh3! 23.gxh3 £d7 24.¥d3 £xh3 25.¦c2 ch4 e l'attacco del
Nero arriva prima)} 22... ch5 23.¢h1 gxh2 24.¥f2 ¥xb5 25.¥xb5 ¤g3+
26.¥xg3 fxg3.

\textbf{20... g3!} se 20... ¥a4? 21.£xa4 ¦xc7 22.¦xc7 £xc7 23.£c2! £d7
24.fxg4 ¤xg4 25.¥xg4 £xg4 26.¤f3

\textbf{21.¤xa8?} si imponeva 21.hxg3 fxg3! 22.¥xg3 ¥h6! 23.¤xa8!
\emph{(23.¤e6 ¥xe6 24.£xe6 ¦g7 e vince)} 23... ch5 24.¥f2 \emph{(24.¥h2
¥e3+ 25.¦f2 £h4 26.¤d3 ¤gf4 27.£e1 ¤xg2! 28.¢xg2 ¦g7+ e vince)} 24...
¤gf4 25.¤d3! ¦g7 26.¤xf4 ¥xf4 27.g4! \emph{(27.¦c7? ¤g3! 28.¦xd7 £h4
29.¦xg7+ ¢xg7 30.¥xa7 ¤xe2+! 31.£xe2 ¥h2+ e vince)} 27... ¥xc1 28.£xc1
¤f4 29.£e3 h5!

\textbf{21... ch5! 22.¢h1} se 22.¥xa7 £h4 23.h3 ¥xh3 24.gxh3 £xh3 25.¦f2
gxf2+ 26.¢xf2 ch4 27.¥f1 £h2+ 28.¤g2 ¦g7 e il Nero vince.

\textbf{22... gxf2 23.¦xf2 ¤g3+! 24.¢g1 £xa8 25.¥c4 a6!} la Donna
rientra in gioco sulla diagonale h7-g1!

\textbf{26.£d3?! £a7 27.b5 axb5 28.¥xb5 ch1! 29.} abbandona.

\bookdiagramplaceholder{assets/diagrams/image77.png}{77}

Una chiusura veramente notevole!

Le altre posizioni che permettono attacchi molto violenti, anche in
mancanza di debolezze posizionali, si hanno con gli arrocchi eterogenei.
In tali casi l'attacco non solo è lecito ma necessario.

\textbf{1. Gli arrocchi eterogenei}

Si dicono \emph{partite con arrocchi eterogenei} quelle in cui un
giocatore arrocca sul lato di Re mentre l'altro su quello di Donna.

Quando ciò si verifica la partita si trasforma in una lotta in cui
tendono a predominare gli elementi tattici. Il motivo dominante della
partita diventa in questi casi l'attacco al re nemico. Per far ciò
entrambi i giocatori spingeranno a fondo i propri pedoni sull'arrocco
cercando di cambiarli con quelli avversari per aprire le linee.

Questa strategia, di norma non possibile, è qui indicata perché le
spinte di pedoni non indeboliscono il proprio arrocco, situato
dall'altra parte della scacchiera.

Ostapenko--Jarzew, URSS 1969 \textbf{(Difesa Siciliana)}

\textbf{1.e4 c5 c}on questa mossa il Nero rompe drasticamente la
simmetria dei pedoni e conferisce alla partita, almeno a grandi linee,
caratteri riconoscibili: il Nero tenderà ad espandersi sul lato di Donna
mentre il Bianco su quello di Re.

\textbf{2.¤f3 ¤c6 3.d4 ¤xd4 4.¤xd4 ¤f6 5.¤c3 d6 6.¥c4 e6 7.¥e3 ¥e7 8.£e2
0-0 9.¥b3 £c7 10.0-0-0} il dado è tratto! Gli arrocchi eterogenei danno
alla partita un motivo strategico dominante: l'attacco diretto al Re. La
posizione richiede la massima energia e una veloce mobilitazione delle
forze.

\textbf{10... a6} il Nero si accinge a spingere in b5

\textbf{11.¦hg1} e il Bianco in g4. I due giocatori si preparano a
lanciare sull'arrocco nemico i propri pedoni.

\textbf{11... b5 12.g4 b4 13.¤xc6} preludio a una combinazione tattica
tesa a recuperare il leggero ritardo nell'attacco ad est rispetto a
quello avversario ad ovest.

\textbf{13... £xc6}

\bookdiagramplaceholder{assets/diagrams/image78.png}{78}

\textbf{14.¤d5! exd5 15.g5 £xe4} se 15... ¤xe4 16.¥xd5.

\textbf{16.gxf6 ¥xf6 17.¥d5} dalla baruffa tattica il Bianco ha ottenuto
l'allungamento della diagonale a2-g8 all'Alfiere campo chiaro e
l'apertura della colonna \emph{g} alla ¦g1. Il Nero ha ottenuto di
rimpiazzare l'Alfiere delle case scure su una diagonale (a1-h8) che
punta direttamente sull'arrocco avversario.

\textbf{17... £a4} speculando sull'intoccabilità della ¦a8. 18.¥xa8
£xa2! e il Bianco, per salvarsi dal matto deve cedere materiale. Per
esempio: 19.¢d2 £xb2 20.¢e1 ¥c3+ 21.¢f1 ¥h3+ 22.¦g2 ¥xg2+ 23.¢xg2 ¦xa8

\textbf{18.£h5} non c'è tempo per la difesa. Occorre contrattaccare a
tutta forza.

\textbf{18... ¥e6} l'inchiodatura del Pf7 impedisce 18... g6 cui
seguirebbe 19.¦xg6+! hxg6 20.£xg6+. Questa mossa contrasta al Bianco la
diagonale a2-g8 e lascia libera la Torre l'importante casa c8.

\textbf{19.¦xg7+!} il Bianco forza i tempi; se 29... ¢xg7 allora
20.£h6+.

\textbf{19... ¥xg7 20.¦g1!} il Bianco minaccia 21.£h6 e poi matto
(£xg7).

\textbf{20... ¦fc8} la lotta si fa drammatica. La ¦a8 è intoccabile per
il matto in c2.

\textbf{21.¦xg7+} un'altra mossa spettacolare. Il Bianco non bada al
materiale pur di aprire l'arrocco avversario. Meno chiara sarebbe stata
21.¥xe4 ¢f8 22.£xh7 ¥xb2+!

\textbf{21... ¢xg7 22.£h6+ ¢g8 23.¥xe4 b3} il Nero riprende il filo
interrotto dell'attacco sull'arrocco, ma non c'è più tempo.

\textbf{24.¥xh7+ ¢h8 25.¥f5+} una mossa precisa il cui scopo diviene
chiaro tra qualche mossa.

\textbf{25... ¢g8 26. £h7+ ¢f8 27.¥h6+ ¢e8 28.£g8+ ¢e7 29.¥g5+ ¢d7
30.£xf7+ s}enza la previdente 25.¥f5 questa mossa non sarebbe stata
possibile. Ora il Re nero è veramente in balia delle onde.

\textbf{30... ¢c6 31.¥xe6 ¢b6 a}lla ricerca di un rifugio per il Re.

\textbf{32.¥e3+ ¢a5 33.¥xc8 ¦xc8 34.£f5+ ¦c5 35.¥xc5 £d5 36.¥b4+ ¢xb4
37.a3+ ¢c4 38.£xb5+ axb5 39.¤xb3+ ¢d3 40.¢d1} abbandona.

In precedenza abbiamo già visto in azione la Variante di cambio del
Gambetto di Donna rifiutato e abbiamo osservato che essa può essere il
preludio a un attacco di minoranza ad ovest.

Nella partita che segue, giocata all'internazionale di Toluca del 1982,
il Bianco rinuncia all'attacco di minoranza per una partita di attacco e
contrattacco con arrocchi eterogenei.

Hulak--Spassky, Toluca 1982 \textbf{(Gambetto di Donna Rifiutato)}

\textbf{1.d4 d5 2.c4} il tentativo di incrementare l'iniziativa con la
prematura spinta in e4 è destinato all'insuccesso: 2.e4 £xe4 3.¤c3 e5! e
il Nero ha risolto i problemi d'apertura.

\textbf{2... e6 3.¤c3 ¤f6 4.¤xd5 exd5 5.¥g5 ¥e7 6.e3 c6 7.£c2 ¤bd7 8.¥d3
0-0 9.¤ge2} decidendo di arroccare lungo questa mossa è la più naturale.
Perde di significato 9.¤f3, utile a rafforzare un arrocco corto ma di
intralcio in caso i arrocchi eterogenei perché ostacola l'avanzata del
pedone ``f''.

\textbf{9... ¦e8 10.h3} taglia definitivamente i ponti con l'arrocco
corto. Il senso dello sviluppo del cavallo in e2 e dell'arrocco corto
poteva essere una strategia di espansione centrale: f2-f3, e3-e4 e,
eventualmente, e4-e5.

\textbf{10... ¤f8 11.0-0-0 a5} Spassky non ci pensa due volte a lanciare
la massa di pedoni sull'arrocco.

\textbf{12.¢b1} Gergadze indica la seguente continuazione contro
l'immediata 12.g4: 12... a4; e ora 13.¤xa4?! £a5 14.¥xf6 (14.¤c3 b6 e
poi ¥b7 e c5) 14... ¥xf6 15.b3 b5.

\textbf{12... b5 13.g4 a4} occorre conquistare la casa a4 prima di
spingere in b4, altrimenti, all'immediata b4, il Bianco gioca ¤a4 e
l'attacco del Nero è frenato.

\textbf{14.¤g3 a3 15.b3 £a5 16.¦hg1 ¢h8} il Re si sposta dalla colonna
ove è posta la Torre.

\bookdiagramplaceholder{assets/diagrams/image79.png}{79}

\textbf{17. ¤ce2 ¥d7 18.¤f5 ¥xf5 19.gxf5} ovviamente non 19.¥xf5 e la
colonna \emph{g} rimane chiusa.

\textbf{19... ¦ac8 20.¤f4 ¤8d7 21.£e2 c5 22.£xc5 ¤xc5 23.¥xf6?} il
Bianco va a caccia di pedoni. Era migliore 23.£e1.

\textbf{23... ¥xf6 24.¤xd5 ¤a4!}

Impedisce il cambio del Cavallo con l'Alfiere e si appresta a sfruttare
la casa c3. Il Cavallo è naturalmente tabù.

\textbf{25.¦c1 ¤c3+ 26.¤xc3?} il Bianco dà troppo peso al materiale.
Occorreva ridare fiato all'attacco giocando 26.¦xc3 ¥xc3 27.f6, con
qualche possibilità. Così si riduce a un gioco passivo.

\textbf{26... ¦xc3 27.¦gd1 £d4} il Nero minaccia 28... ¦xb3+ e matto in
poche mosse.

\textbf{28.¥c2 T¤xe3 29.£d2} la Torre nera è intoccabile: 29.fxe3 £c3!

\textbf{29... ¦c3 30.¦e1 ¦xe1 31.£xe1 h6 32.¦d1 ¢h7 33.£e2 ¦xh3 34.£e1
£c5 35.¢c1}

\bookdiagramplaceholder{assets/diagrams/image80.png}{80}

\textbf{35... ¦xb3 36.axb3 a2 37.} abbandona.

\textbf{2. L'indebolimento dell'arrocco}

\textbf{a) Indebolimento per spinta di pedoni}

Un motivo che giustifica un attacco, eventualmente appoggiato da pedoni,
si ha quando l'arrocco avversario si indebolisce a causa di una spinta
di pedone.

\bookdiagramplaceholder{assets/diagrams/image81.png}{81}

Questo è il più comune indebolimento dell'arrocco. La spinta in h6 può
essere necessaria per scacciare un Alfiere da g5 o per concedere una
casa di fuga al Re nero ma è pur sempre un indebolimento.

Occorre evitare di fare questa spinta prima che l'avversario abbia
arroccato sullo stesso lato, altrimenti egli arroccherà sull'altra ala e
spingerà i propri pedoni sull'arrocco nemico. Questa strategia è
particolarmente efficace perché il pedone h6 fa da punto di rottura.

In altre parole la strategia del Bianco sarà di spingere un pedone in g5
e aprire la colonna g.

Schema di attacco contro l'indebolimento causato dalla spinta di un
passo del pedone di Torre.

\bookdiagramplaceholder{assets/diagrams/image82.png}{82}

\bookdiagramplaceholder{assets/diagrams/image83.png}{83}

Un pedone davanti all'arrocco può essere spinto anche per porre
l'Alfiere in fianchetto (diagramma 83). In questo caso l'arrocco,
potendo contare anche su un pezzo, ne esce addirittura rafforzato a
patto però che l'Alfiere in fianchetto non venga cambiato.

La struttura dei pedoni corretta, in caso di sviluppo dell'Alfiere in
fianchetto, è quella mostrata in diagramma; occorre evitare di spingere
il pedone ``e'' per non bucare ulteriormente la posizione sulle case
nere (al buco in h6 si aggiungerebbe quello in f6 e si indebolirebbe la
diagonale nera a3-f8.

Un piano di attacco del Bianco contro tale struttura è, arrocco lungo,
indebolire l'arrocco avversario mediante il cambio degli Alfieri su casa
scura (¥e3, £d2 e poi ¥h6) e di aprire la colonna ``h'' mediante la
spinta h2-h4-h5-hxg6. Anche in questo caso si vede bene come la spinta
del pedone davanti all'arrocco rende più facile l'apertura delle
colonne.

\bookdiagramplaceholder{assets/diagrams/image84.png}{84}

Il diagramma mostra lo schema di attacco contro un arrocco con
fianchetto.

La partita che segue è una dimostrazione pratica di come condurre
l'attacco.

Karpov--Gik, Mosca, Campionato studentesco 1968 \textbf{(Difesa
Siciliana)}

\textbf{1.e4 c5 2.¤f3 d6 3.d4 ¤xd4 4.¤xd4 ¤f6 5.¤c3 g6} la variante del
Dragone, una delle più complicate e taglienti varianti della Siciliana.

\textbf{6.¥e3 ¥g7} a 6... ¤g4? seguirebbe 7.¥b5+.

\textbf{7.f3 0-0} il Bianco previene il salto di Cavallo in g4.

\textbf{8.¥c4 ¤c6 9.£d2 £a5 10.0-0-0} con gli arrocchi eterogenei si
delinea lo sviluppo della partita. Il bianco pingerà in h5 per aprire la
colonna h, il Nero attaccherà sul lato di Donna cercando di valorizzare
l'Alfiere in fianchetto.

\textbf{10... ¥d7 11.h4 ¤e5} contro l'attacco ad est il Nero
contrattacca sulla colonna ``c''.

\textbf{12.¥b3 ¦fc8 13.h5 ¤xh5 14.¥h6 ¥xh6 15.£xh6 ¦xc3} Karpov scrive:
\emph{``È un tipico sacrificio di qualità in situazioni simili. Da una
parte il Nero sfugge per sempre dalla minaccia dell'assalto del Cavallo
bianco nella casa d5, e dall'altra parte spera di arrivare più presto al
Re bianco''.}

\textbf{16.bxc3 £xc3 17.¤e2 £c5 18.g4 ¤f6 19.g5 ch5 20.¦xh5} il Bianco
ha fretta. Abbiamo visto anche in altre occasioni come in partite i
attacco e contrattacco con arrocchi eterogenei i tempi possano contare
più del materiale. A 20.¤g3 il Nero avrebbe risposto con 20... ¥g4! e la
Donna bianca rimane fuori gioco.

\textbf{20... gxh5 21.¦h1 £e3+ 22.¢b1} una mossa precisa. Su 22. ¢b2? il
Nero ha almeno la patta: 22... ¤d3+ 23.¤xd3 (23.¢b1?? £xf3! e vince il
Nero!) 23... £xe2+ 24.¢a1 £xd3 e il Nero ha almeno lo scacco perpetuo.

\textbf{22... £xf3} 22... £xe2? 23.£xh5 e6 24.£xh7+ ¢f8 25.£h8+ ¢e7
26.£f6+ ¢e8 27.¦h8 matto.

\textbf{23.¦xh5 e6 24.g6 ¤xg6 25.£xh7+ ¢f8}

\bookdiagramplaceholder{assets/diagrams/image85.png}{85}

\textbf{26.¦f5!!} una mossa molto bella. Ecco il commento di Evgeny Gik
a questa mossa: \emph{``Questa mossa di Torre è stata per me come un
fulmine a cielo sereno! Questa elegante idea geometrica ha deciso subito
il risultato della partita. Le due linee--la diagonale a2-g8 e la
colonna ``f'' si sono incrociate nel punto critico f7. Minaccia £xf7 con
scacco matto, nello stesso tempo la Torre bianca appoggia la Donna sulla
colonna f, mentre l'Alfiere la appoggerebbe in caso di 26... exf5, sulla
diagonale. Il Nero non ha altra via d'uscita che sacrificare la
Donna.''}

\textbf{26... £xb3+ 27.axb3 exf5 28.¤f4 ¦d8 29.£h6+} Karpov: \emph{``È
l'ultima piccola sottigliezza: il Bianco vuole prendere il pedone g6 con
scacco''.}

\textbf{29... ¢e8 30.¤xg6 fxg6 31.£xg6+ ¢e7 32.£g5+} Karpov:
\emph{``Bisogna essere precisi fino alla fine. Dopo 32.exf5 ¦f8 il Nero
avrebbe potuto ancora resistere. Invece adesso non c'è alcuna difesa
all'irruente avanzata del pedone f~''}.

\textbf{32... ¢e8 33.exf5 ¦c8 34.£g8+ ¢e7 35.£g7+ ¢d8 36.f6} il Nero
abbandona.

\textbf{b) Indebolimento per rimozione di pezzi.}

La spinta di un pedone davanti all'arrocco produce una debolezza
duratura ma l'arrocco può essere debole anche perché non protetto a
sufficienza dai pezzi.

La partita tra Marshall e Wolf è un bell'esempio di sfruttamento di un
arrocco debole. Il campione americano Frank J. Marshall (1877-1944)
costringe l'avversario a rimuovere l'ultimo importante difensore e poi
scatena la furia devastatrice dei pezzi bianchi sul Re nemico.

\bookdiagramplaceholder{assets/diagrams/image86.png}{86}

Di solito la struttura dei pedoni suggerisce la disposizione dei pezzi e
la direzione di attacco. In questa posizione sono invece indicativi gli
Alfieri di entrambi i colori che puntano direttamente sull'arrocco. I
Cavalli fungono da cani da gua¢dia ma il Bianco, che ha il tratto,
riesce ad allontanare il buon custode con la mossa:

\textbf{16.¤e4!} il Nero non può giocare 16... ¥e7 per il doppio attacco
in f6 che gli rovinerebbe la struttura dei pedoni, né può chiudere la
diagonale a2-g8 con 16... ¤xe4 17.¥xe4 f5 perché dopo 18.¥xc6 ¥xc6
19.¤d4 perderebbe un pezzo; devi quindi spostare il difensore.

\textbf{16... ¤d5 17.¤eg5 g6} necessaria non potendosi giocare 17... h6
18.£c2 g6 19.¤xe6!

\textbf{18.¤xh7! ¢xh7 19.¤g5+ ¢g8} in caso di 19... ¢h6 segue 20.£g4.

\textbf{20.£h5!!} una mossa ad effetto. La Donna è imprendibile per via
del matto: 20... gxh5 21.¥h7 matto.

\textbf{20... f6 21.¥xg6 ¦d7} 21... £a7 22.¦fd1.

\textbf{22.¤xe6 ¦h7 23.¥xh7+ £xh7 24. £xh7+ ¢xh7 25.¤xf8+ ¥xf8 26.¦fd1
cce7 27.e4 ¤b6 28.¦c7 ¢g8 29.¥xf6 ¤g6 30.¦d8} il Nero abbandona.

\textbf{c) Il Re al centro}

Abbiamo avuto modo di vedere come la strategia cambia in modo radicale
al variare della posizione del Re. In caso di sbilanciamento di un
colore su di un'ala, l'altro potrà sempre arroccare sull'ala opposta.
Per questa ragione a volte può risultare vantaggioso ritardare l'arrocco
per effettuarlo quando la partita ha assunto caratteristiche più
accentuate.

Naturalmente questa strategia comporta un certo rischio, dovuto in parte
al ritardo nel completamento dello sviluppo e, in parte, alla
possibilità di un'improvvisa apertura del centro con la conseguente
impossibilità di arroccare. Il Re verrebbe così a trovarsi al centro
della tempesta!

\emph{Partita Spagnola} \textbf{1.e4 e5 2.¤f3 ¤c6 3.¥b5 a6 4.¥a4 ¤f6
5.d3} più comune è 5.00. Questo modo di giocare la Spagnola era tipico
di Steinitz.

\textbf{5... d6 6.c3 ¥e7 7.h3 0-0 8.£e2 ¤e8 9.g4 b5 10.¥c2 ¥b7 11.¤bd2
£d7 12.¤f1 ¤d8 13.¤e3 ¤e6 14.¤f5} il Bianco può intraprendere le
operazioni sul lato di Re mentre il Nero non sa ancora dove dovrà
contrattaccare.

\textbf{14... g6 15.¤xe7+ £xe7 16.¥e3} da quale parte arroccherò? Dalla
parte opposta a quella che pensi, qualsiasi cosa pensi, sembra dire il
Bianco all'avversario.

\textbf{16... ¤8g7 17.0-0-0} dà infine l'indirizzo di casa.

Questo inizio di partita fu giocato da Steinitz contro Blackburne a
Londra nel 1876.

C'è poi il caso in cui la rinuncia all'arrocco non è un tatticismo
provvisorio ma una necessità strategica; per esempio, quando un lato si
è indebolito e nell'altro è in corso una schermaglia, il centro può
essere il rifugio più sicuro per il Re.

\emph{Difesa Francese:} \textbf{1.e4 e6} la Francese è una difesa in cui
il Re al centro (anche Bianco) non è affatto un'eccezione. Questa è una
delle tante varianti.

2.d4 d5 3.¤c3 ¥b4 4.e5 c5 5.a3 ¥xc3+ 6.bxc3 ¤e7 7.£g4 ¤xd4 8.£xg7 ¦g8
9.£xh7 £c7 10.¤e2 ¤bc6 11.f4 ¥d7 12.£d3 £xc3 13.¤xc3 a6 14.¤e2 ¤f5
15.¦b1 ¦c8

È palese che in questa posizione la decisione di non arroccare è la
migliore soluzione ma ci sono anche aperture in cui uno dei due colori
si priva in modo volontario dell'arrocco fin dalle prime mosse, senza
che la teoria sentenzi essere svantaggioso.

Vediamo:

\emph{Gambetto di Donna accettato:} \textbf{1.d4 d5 2.c4 £xc4 3.¤f3 ¤f6
4.e3 e6 5.¥xc4 c5 6.0-0 a6 7.£xc5 ¥xc5 8.£xd8+ ¢xd8} il Re Nero si porrà
successivamente in e7, casa che, in questa variante, sembra
sufficientemente sicura.

\emph{Vecchia Indiana:} \textbf{1.d4 ¤f6 2.c4 d6 3.¤c3 e5} anche qui la
continuazione 4.£xe5 £xe5 5.£xd8 non sembra dare al Bianco particolari
vantaggi. Tant'è che molti preferiscono non cambiare le Donne e giocare:
4.d5.

\emph{Gambetto di Re:} \textbf{1.e4 e5 2.f4 exf4 3.¥c4} senza
preoccuparsi dello scacco in h4.

\textbf{3... £h4+ 4.¢f1} la pratica ha dimostrato che, dopo questa
mossa, i maggiori problemi deve affrontarli il Nero, tant'è che la
continuazione consigliata dalla teoria in risposta a 3.¥c4 è 3... ¤f6.
Con 3... £h4+ 4.¢f1 Bobby Fischer sconfisse il connazionale Grande
Maestro Larry Evans in 36 mosse nel Campionato americano del 1963

In linea di massima, però, il Re non arroccato costituisce un grave
svantaggio.

Vediamo che cosa capita di solito a chi ha commesso l'imprevidenza di
lasciare il Re al centro.

Corvi--Leoncini, S. Benedetto del Tronto 1988 \textbf{(Difesa
Siciliana)}

\textbf{1.e4 c5 2.d4 ¤xd4 3.c3 d5} rientrando nella variante Alapin
(1.e4 c5 2.c3) che di norma gioco col Bianco. Per accettare il gambetto
occorre una approfondita conoscenza teorica e anche in questo caso si
rischia di rimanere invischiati in qualche novità teorica.

\textbf{4.exd5 £xd5 5.¤f3 ¤c6 6.¤xd4 e5} una continuazione molto tattica
sulla cui correttezza è lecito nutrire qualche dubbio.

Nel primo torneo in cui, contro la Siciliana, decisi di passare a 2.c3
(prima giocavo 2.¤f3) il maestro pisano Pierluigi Beggi mi sorprese con
questa mossa (questo la dice lunga su come mi fossi preparato!). Ne
rimasi così colpito da aggiungerla al mio repertorio.

\textbf{7.£xe5?!} la stessa mossa che giocai con Beggi!
\emph{``Possibile che il Nero possa compiere la spinta liberatoria in e5
dopo appena sei mosse?~''} mi domandai, e decisi che tale sfrontatezza
andava punita subito! Visto come sono andate le cose, immagino che Marco
Corvi sia stato vittima di un ragionamento simile. Ora il Nero conquista
l'iniziativa e si appresta ad attaccare il Re al centro. La
continuazione canonica è invece 7.¤c3 ¥b4 8.¥d2 ¥xc3 9.¥xc3 e4 10.¤e5
¤xe5 11.£xe5 ¤e7 12.£c2 e il Bianco rimane col lieve, ma noioso,
vantaggio della coppia degli Alfieri.

\textbf{7... £xd1+ 8.¢xd1 ¥g4} ricordo ancora l'espressione disgustata
del mio avversario. Aver deciso di giocare all'attacco e trovarsi in
difesa dopo poche mosse, per giunta col Bianco, lo mise di sicuro in
grave svantaggio psicologico.

\textbf{9.¥f4 ¥c5} nella partita contro di me Beggi preferì continuare
con 9... ¤ge7 e mi vengono ancora oggi i brividi a ripensarci. Ecco come
riuscii a vincere quella partita partendo da una posizione così poco
attraente: 10. ¥b5 0-0-0+ 11.¤bd2 ¤d5 (Probabilmente era migliore 11...
¤g6) 12.¥g5 f6 13.exf6 gxf6 14.¥h4 ¥xf3+ 15.gxf3 ¥h6 16.¥xc6 ¤b4 17.¥e4
¦xd2+ 18.¢e1 ¦e8 19.¦c1+ (Il Bianco conquista l'iniziativa e colpisce in
contropiede!) 19... ¢b8 20.¥g3+ ¢a8 21.¢f1 ¦d7 22.¦c4 ¤xa2 23.¢g2 ¥f8
24.¦a1 ¤b4 25.¦ac1 ¤c6 26.¦xc6!! il Nero abbandona (Leoncini--Beggi,
Gaeta 1987).

\textbf{10.¥g3} la mossa migliore è 10.¤bd2!, come mi giocò il Candidato
Maestro torinese Roberto Pietrocola quattro anni dopo a Imperia, mossa
che difende indirettamente il Pf2. A 10... ¥xf2 seguirebbe infatti
11.¤e4! e il Bianco riconquista l'iniziativa con gli interessi.

\textbf{10... ¤ge7 11.¤bd2 0-0} una valida alternativa a 0-0-0.
L'arrocco lungo, se è vero che mette in gioco prima la Torre di Donna,
può risultare fatale al Re in caso di controgioco nero, come mostra la
mia partita con Beggi.

\textbf{12.¥e2 ¦fd8} la Torre giusta, l'altra è destinata ad andare in
c8.

\textbf{13.¦c1 ¥b6} sono state effettuate solo tredici mosse dall'inizio
della partita e la posizione del Bianco è già disperata. In queste
condizioni lo sconforto è così grande da indurre, esso stesso,
all'errore.

\textbf{14.¢c2? ¥f5+ 15.¢b3}

\bookdiagramplaceholder{assets/diagrams/image87.png}{87}

\textbf{15... ¦xd2 16.¤xd2 ¤d4+ 17.¢a3 ¤xe2} il guadagno di materiale
non diminuisce affatto l'attacco. Il Bianco potrebbe abbandonare.

\textbf{18.T¤d1 ¥c2 19.T£e1 ¤d4 20.b3 ¥f5 21.¦d1 ¤c2+ 22.¢b2 ¥d4+ 23.¢c1
¤b4 24.¤c4 ¤xa2+ 25.¢d2 ¦c8 26.¢e1 ``}Non c'è pace tra gli ulivi''.

26... ¥c5 27.¤d6 ¥b4+ 28.¢e2 ¥g4+ 29.¢e3 ¦c2 30.f3 ¥e6 31.¤xb7 h6
32.¦d8+ ¢h7 33.¤d6 ¤d5+ 34.¢d3 ¦d2+

\bookdiagramplaceholder{assets/diagrams/image88.png}{88}

e matto alla prossima. Graziosi i matti con mosse di Cavallo: 35.¢c4 ¤c7
matto; 35.¢e4 ¤ac3 matto.

\textbf{35.} il Bianco abbandona.
